\section*{第 \ref{ch:b3:hamiltonian-mechanics} 章——哈密顿力学}

\begin{solution}{exo:b3:hamiltonian-mechanics:1}
(a) $p = m\dot x$，$H = p\dot x - L = p^2/2m + \tfrac12 kx^2 = E_k +
E_p$。(b) $\dot x = p/m$，$\dot p = -kx$，因此 $\ddot x = -(k/m)x$。
(c) $H = E$ 为椭圆 $x^2/(2E/k) + p^2/(2mE) = 1$：半轴分别为
$\sqrt{2E/k}$ 和 $\sqrt{2mE}$。(d) 在最右侧点（$x > 0$，
$p = 0$）处，$\dot p = -kx < 0$：点向下运动——顺时针，
始终如此。
\end{solution}

\begin{solution}{exo:b3:hamiltonian-mechanics:2}
(a) $p = mR^2\dot\theta$；$H = p^2/2mR^2 - \tfrac12 m\omega^2R^2
\sin^2\theta - mgR\cos\theta = p^2/2mR^2 + U_{\text{eff}}(\theta)$。
(b) 守恒（$H$ 不显含 $t$），但它是 $h$，而非
机械能：电动机做的功未计入其中。(c) 等值线
为 $p = \pm\sqrt{2mR^2(H - U_{\text{eff}})}$：慢速
转动时，在 $(0, 0)$ 附近为卵形曲线，并过
鞍点 $\theta = \pi$ 处的分界线；快速转动时，在
$(\pm\theta_{\text{平衡}}, 0)$ 附近为两族卵形曲线。(d) 快速情形：中心位于
$\pm\theta_{\text{平衡}}$；鞍点位于 $\theta = 0$ \emph{与} $\theta =
\pi$；过 $\theta = 0$ 的是一条八字形分界线，
包围两个中心（小振动在底部跳跃），而
过 $\pi$ 的是外侧分界线，珠子在其外
绕过最高点运动。
\end{solution}

\begin{solution}{exo:b3:hamiltonian-mechanics:3}
(a) 一次飞行对应弧 $p = \pm\sqrt{2m^2g(z_{\max} - z)}$——一条
侧卧的抛物线，自 $(0, +p_0)$ 上升至
$(z_{\max}, 0)$，再下降至 $(0, -p_0)$。(b) 反弹将
$(0, -p_0)$ 映为 $(0, +p_0)$：一段竖直线段，使轨道闭合。(c)
$\oint p\,\dd z = 2\int_0^{z_{\max}}m\sqrt{2g(z_{\max} - z)}\,\dd z =
\tfrac43 m\sqrt{2g}\,z_{\max}^{3/2}$，而 $I$ 为其除以 $2\pi$。(d)
$z_n = \big[3nh/(4m\sqrt{2g})\big]^{2/3}$：对中子，$z_1 =
(\num{6.7e-8})^{2/3} = \qty{17}{\micro m}$——量级正确：
格勒诺布尔实验测得最低引力量子态位于
约 \qty{15}{\micro m}（精确处理采用真实波函数，
给出 \qty{14}{\micro m}）。
\end{solution}

\begin{solution}{exo:b3:hamiltonian-mechanics:4}
(a) $\{x^2, p\} = 2x$。(b) $\{xp, H\} = p^2/m - x\,E_p'(x) = 2E_k -
x\,E_p'$：对束缚运动，$\dd(xp)/\dd t$ 的时间平均为零，
故 $\langle 2E_k\rangle = \langle x\,E_p'\rangle$——即
维里定理（对 $E_p \propto x^2$：$\langle E_k\rangle = \langle
E_p\rangle$；对 $\propto -1/r$：$2\langle E_k\rangle = -\langle
E_p\rangle$）。(c) $\{L_z, x\} = y$，$\{L_z, p_x\} = p_y$：$L_z$
同时生成位置与动量的转动。(d) 取 $h
= H$ 的雅可比恒等式：$\{\{f, g\}, H\} = -\{\{g, H\}, f\} - \{\{H, f\}, g\} = 0$。
\end{solution}

\begin{solution}{exo:b3:hamiltonian-mechanics:5}
(a) $E_{\text{分界}} = mg\ell$（倒立静止位置的
能量）；在最低点 $|p| = 2m\ell^2\omega_0$ 取最大值。(b) 由
$p^2/2m\ell^2 = mg\ell(1 + \cos\theta) = 2mg\ell\cos^2(\theta/2)$ 得出。(c)
$\dot\theta = 2\omega_0\cos(\theta/2)$ 可分离变量：$\omega_0t =
\ln\tan(\theta/4 + \pi/4)$，当 $\theta \to \pi$ 时发散：
趋近最高点但永不到达。(d) 恰在分界线下方，运动
紧随分界线：摆蠕动进入鞍点的邻域，
停留良久——即对数项所示——最终落回；每个
鞍点都使轨迹按对数方式变慢，这正是一根杆
未能完美平衡的杆在倒下前显得犹豫的原因。
\end{solution}

\begin{solution}{exo:b3:hamiltonian-mechanics:6}
(a) $\dot{\vect r} = (\vect p - q\vect A)/m$ 及 $\dot p_x =
-\partial H/\partial x = (q/m)(\vect p - q\vect A)\cdot\partial_x\vect
A$ 等；消去 $\vect p$ 即得 $m\ddot x = qB\dot y$，
$m\ddot y = -qB\dot x$。(b) $H = (\vect p - q\vect A)^2/2m = \tfrac12
m\vect v^{\,2}$：仅为动能——为常量，因为
磁力与 $\vect v$ 垂直。(c) 由 $m\dot x = p_x +
\tfrac12 qBy$：$\pi_x = p_x - \tfrac12 qBy = m\dot x - qBy$，其
守恒即为运动方程的第一式（同理 $\pi_y$）；在泊松
括号中仅两项保留：$\{p_x, \tfrac12 qBx\} = -\tfrac12
qB$ 与 $\{-\tfrac12 qBy, p_y\} = -\tfrac12 qB$，合计 $-qB$。（$\pi_x/qB$ 与 $\pi_y/qB$
在相差符号之内，即圆周的引导中心
的坐标。）(d) 由 \cref{exo:b3:hamiltonian-mechanics:4}(d) 可知，
两个守恒量的泊松括号仍守恒——此处为
常数 $-qB$，它平凡守恒，未带来新信息：没有第三个
守恒量出现。
\end{solution}

\begin{solution}{exo:b3:hamiltonian-mechanics:7}
(a) 映射 $(q, p) \mapsto (q + pt/m, p)$ 是剪切：底与
高不变，面积亦然（行列式为 $1$）。
(b) 在变量 $(x\sqrt{m\omega}, p/\sqrt{m\omega})$ 下，流为
以角速率 $\omega$ 的刚性转动；转动保持面积，
且变量替换的行列式为 $1$。(c) 速度场
$(v, -\omega^2x - \gamma v)$ 的散度为 $-\gamma$：任一面积以
$\eu^{-\gamma t}$ 收缩，螺旋趋向原点——
哈密顿流中不可能出现。(d) 流入被忽略的
自由度：空气分子与金属丝的声子，其相空间
体积至少同等地增长——这正是热力学第二定律的种子。
\end{solution}

\begin{solution}{exo:b3:hamiltonian-mechanics:8}
(a) $\dot r = p_r/m$，$\dot\varphi = p_\varphi/mr^2$，$\dot p_r =
p_\varphi^2/mr^3 - E_p'(r)$，$\dot p_\varphi = 0$。(b) $H$ 不
含 $\varphi$：$\{p_\varphi, H\} = -\partial H/\partial\varphi =
0$；角动量守恒。(c) $H_{\text{辐射}} = p_r^2/2m +
U_{\text{eff}}(r)$，其中 $U_{\text{eff}} = p_\varphi^2/2mr^2 + E_p(r)$。
(d) $U_{\text{eff}}' = 0$：$r_{\text{c}} = p_\varphi^2/mk$，且 $E =
U_{\text{eff}}(r_{\text{c}}) = -mk^2/2p_\varphi^2$.
\end{solution}

\begin{solution}{exo:b3:hamiltonian-mechanics:9}
(a) $\{L_x, L_y\} = \{yp_z - zp_y,\ zp_x - xp_z\}$：仅有的
非零正则括号给出 $yp_x\{p_z, z\} + xp_y\{z,
p_z\} = -yp_x + xp_y = L_z$。(b) 作循环轮换 $x \to y \to z \to
x$ 即得另外两式。(c) $\{L^2, L_z\} = 2L_x\{L_x, L_z\} +
2L_y\{L_y, L_z\} = -2L_xL_y + 2L_yL_x = 0$。(d) 对 $H = L^2/2J$，
每个 $\{L_i, H\} = 0$：$\vect L$ 保持不变。若 $H = \sum L_i^2/2J_i$
且 $J_i$ 互不相等：$\{L_x, H\} = L_yL_z(1/J_y - 1/J_z) \neq 0$——
仅 $L^2$ 与 $H$ 保持守恒，刚体发生翻滚（本科二年级
分册刚体一章中的网球拍定理即在此体现）。
\end{solution}

\begin{solution}{exo:b3:hamiltonian-mechanics:10}
(a) $\oint p\,\dd q = 2\pi I = 2\pi E/\omega = nh$ 可得 $E_n =
n\hbar\omega$。(b) 在 $|p|$ 恒定时往返运动：$2pa = nh$，$p =
nh/2a$，$E = p^2/2m = n^2h^2/8ma^2$——恰为
本科二年级分册中的无限深势阱。(c) $\omega = \sqrt{k/\mu} = \qty{4.1e14}{rad/s}$，
$\hbar\omega = \qty{4.3e-20}{J} = \qty{0.27}{eV}$, $\lambda = hc/
\hbar\omega = \qty{4.6}{\micro m}$：中红外（CO 基频
带，用于追踪星际气体）。(d) 否：能级之差
不受共同的一半影响。零点能体现于
依赖绝对能级的比较——离解能的同位素
位移（H$_2$ 与 D$_2$ 对比），或
晶体中在绝对零度下仍存在的振动。
\end{solution}

\begin{solution}{exo:b3:hamiltonian-mechanics:11}
(a) $E = p^2/2m + \tfrac12 k(t)x^2$，故沿运动有 $\dot E =
\partial E/\partial t = \tfrac12\dot kx^2$。(b) 在一个周期内，
$k$ 近似不变时，$\langle\tfrac12 kx^2\rangle = E/2$，故
$\langle\dot E\rangle = \dot k\,E/2k$。(c) $\dd\ln E = \tfrac12\dd\ln
k$；由 $\omega = \sqrt{k/m}$ 可得 $\dd\ln\omega = \tfrac12\dd\ln k$
同样成立：$E/\omega$ 不变。(d) $\omega \propto \ell^{-1/2}$ 增大
为原来的 $\sqrt2$ 倍，故 $E$ 增大 $\sqrt2 \approx 1.41$ 倍。代入 $E = \tfrac12
mg\ell\theta_0^2$，$\theta_0 \propto \ell^{-3/4}$：$\theta_0' =
5^\circ \times 2^{3/4} = 8.4^\circ$。能量由
拉绳者提供：张力平均超过 $mg\cos\theta$
（离心项），因此提拉对摆做正的净功。
\end{solution}

\begin{solution}{exo:b3:hamiltonian-mechanics:12}
(a) $\oint p\,\dd\theta = 2\int_{-\pi}^{\pi}2m\ell^2\omega_0
\cos(\theta/2)\,\dd\theta = 4m\ell^2\omega_0\big[2\sin(\theta/2)
\big]_{-\pi}^{\pi} = 16m\ell^2\omega_0$. (b) $\omega_0 =
\qty{9.9}{rad/s}$: $16 \times \num{e-3} \times \num{e-2} \times 9.9 =
\qty{1.6e-3}{J.s}$，即 $N \approx \num{2.4e30}$ 个态：
实验室单摆的量子颗粒性，
比任何可观测量低三十个数量级。(c) $16 \times \num{e-26} \times
\num{e-20} \times \num{e14} = \qty{1.6e-31}{J.s}$，即 $N \approx
240$：分子摆动仅容纳几百个量子态，
其低能级可由光谱逐一分辨。(d) 界限
在于运动的相空间面积
仅为 $h$ 的适度倍数：分子及更小尺度为量子的；凡可见之物皆为经典的。
\end{solution}

\begin{solution}{pb:b3:hamiltonian-mechanics:1}
\textbf{1.} $m_{\text{p}}/m_{\text{e}} = 1836$：质子运动小 1836
倍；$H$ 如题所述，位于轨道平面内。
\textbf{2.} $\dot r = p_r/m_{\text{e}}$, $\dot\varphi =
p_\varphi/m_{\text{e}}r^2$, $\dot p_r = p_\varphi^2/m_{\text{e}}r^3 -
k/r^2$, $\dot p_\varphi = 0$.
\textbf{3.} $\varphi$ 不出现于 $H$ 中：$p_\varphi$，角
动量守恒。
\textbf{4.} $U_{\text{有效}} = p_\varphi^2/2m_{\text{e}}r^2 - k/r$：
小 $r$ 处为排斥壁，大 $r$ 处为库仑尾部，
二者之间有一个极小值。
\textbf{5.} 由 $U_{\text{有效}}' = 0$ 得 $r_{\text{c}} =
p_\varphi^2/m_{\text{e}}k$.
\textbf{6.} $E = k/2r_{\text{c}} - k/r_{\text{c}} = -k/2r_{\text{c}} =
-m_{\text{e}}k^2/2p_\varphi^2$；$E_p = -k/r_{\text{c}} = 2E$：
任意 $1/r$ 型圆轨道的位力比。
\textbf{7.} $m_{\text{e}}v^2/r = k/r^2$: $v = \sqrt{k/m_{\text{e}}r}$,
$f_{\text{轨道}} = v/2\pi r = (1/2\pi)\sqrt{k/m_{\text{e}}r^3}$。
\textbf{8.} $k = \qty{2.30e-28}{J.m}$: $v = \qty{2.2e6}{m/s}$ ($v/c =
0.0075$），$f_{\text{轨道}} = \qty{7.2e15}{Hz}$，$E = -k/2r =
\qty{-2.3e-18}{J} = \qty{-14}{eV}$.
\textbf{9.} $p_\varphi = m_{\text{e}}vr = \qty{1.0e-34}{J.s} \approx
\hbar$：原子轨道携带的角动量量级为 $\hbar$
——普朗克常量内蕴于原子尺度之中。
\textbf{10.} 每个原子应在 $\sim\qty{e-11}{s}$ 内坍缩，其
光连续扫向更短的波长。观测结果：原子
永恒存在，并发出尖锐、固定的谱线。
\textbf{11.} $1/\lambda = R_H(1/n^2 - 1/n'^2)$ 将每一观测频率
写成两项之\emph{差}：能量在
固定能级 $-hcR_H/n^2$ 之间交换。
\textbf{12.} 锁定于普适值的量，在
原子受缓慢扰动（场、碰撞、缓慢变化）时不得漂移；一个绝热
不变量恰是在缓慢
扰动下保持不变的量——因此将作用量量子化。
\textbf{13.} $p_\varphi$ 在轨道上为常量：$\oint
p_\varphi\,\dd\varphi = 2\pi p_\varphi = nh$，即 $p_\varphi =
n\hbar$.
\textbf{14.} $r_n = (n\hbar)^2/m_{\text{e}}k = n^2a_0$, $a_0 =
\hbar^2/m_{\text{e}}k = \qty{5.29e-11}{m} = \qty{52.9}{pm}$.
\textbf{15.} $E_n = -m_{\text{e}}k^2/2n^2\hbar^2 = -E_{\text{I}}/n^2$,
$E_{\text{I}} = m_{\text{e}}k^2/2\hbar^2 = \qty{2.18e-18}{J} =
\qty{13.6}{eV}$.
\textbf{16.} $v_1 = k/\hbar = \qty{2.19e6}{m/s}$; $v_1/c = k/\hbar c =
1/137$：精细结构常数 $\alpha$，衡量
原子的非相对论性程度。
\textbf{17.} $h\nu = E_{n'} - E_n$ 给出 $1/\lambda =
(E_{\text{I}}/hc)(1/n^2 - 1/n'^2)$: $R_H = E_{\text{I}}/hc =
\qty{1.10e7}{m^{-1}}$，与实测的 \qty{1.097e7}{m^{-1}} 在
所用常数的精度范围内相符。
\textbf{18.} $2 \to 1$：\qty{122}{nm}（远紫外，莱曼
$\alpha$）；$3 \to 2$：\qty{656}{nm}，使发射星云
呈色的可见红线；$\infty \to 2$：\qty{365}{nm}，巴耳末系的近紫外
边界。
\textbf{19.} $k \to 2k$：半径缩小 $2$ 倍，$E_{\text{I}} \to
4 \times \qty{13.6}{eV} = \qty{54.4}{eV}$；$2 \to 1$ 谱线位于
$\qty{122}{nm}/4 \approx \qty{30}{nm}$，处于极紫外区。
\textbf{20.} $f_{\text{轨道}}(n) = 2E_{\text{I}}/hn^3$；
$\nu_{n \to n-1} = (E_{\text{I}}/h)(2n - 1)/n^2(n-1)^2$。比值
$\nu/f_{\text{轨道}} = n^3(2n-1)/2n^2(n-1)^2$：$3$ 对应 $n = 2$，$1.17$
对应 $n = 10$，$1.015$ 对应 $n = 100$。
\textbf{21.} 在量子数很大的极限下，量子
预言必须并入经典结果——此处即辐射
频率并入轨道频率。
\textbf{22.} $n = 0$ 意味着 $p_\varphi = 0$：一个
“圆轨道”，半径为零，电子落在质子上——不存在这样的运动。
\textbf{23.} 实测氢原子的基态具有\emph{零}
角动量（而非 $\hbar$），且有若干不同的态共享
同一 $n$（化学中的 $s$、$p$、$d$ 形状）——两者
对单个圆轨道而言均不可能。
\textbf{24.} $a_0 \propto 1/m$: $a_0^\mu = \qty{256}{fm}$;
$E_{\text{I}}^\mu = 207 \times \qty{13.6}{eV} = \qty{2.8}{keV}$
（X 射线）。μ子的轨道半径缩小约两百倍——部分
\emph{深入}重元素的核电荷之内——因此其能级
可用来测量核半径（并使质子半径的现代谜题
更为尖锐）。
\textbf{25.} 一个绝热不变量取为 $nh$，即可给出 $a_0 =
\qty{52.9}{pm}$, $E_{\text{I}} = \qty{13.6}{eV}$, $R_H =
\qty{1.10e7}{m^{-1}}$——这些数值，在
\cref{ch:b3:hydrogen-atom}中的完整量子理论将保持这些数值不变，同时取代
产生它们的轨道。
\end{solution}
